\section{Discussion}
\label{sec:discussion}
\textbf{Measure before you automate.} The fixed-action matrix contradicts the default runbook: restarting
the component that looks broken is the right answer for only two of five training faults, and a
wrong action does not degrade gracefully, it prolongs the outage to the timeout. Any remediation
automation, human-written or learned, should be validated against measured outcomes.

\textbf{Learned policies generalise through symptoms.} Because the bandit sees symptoms rather than
fault names, it transfers to a fault it never saw when that fault presents like a known one but
needs a different fix than its look-alike's runbook (the partition). It does not compose actions:
the compound fault needs a two-step plan that a one-step learner cannot represent. The learner also
inherits ambiguity from its inputs: pool timeouts make a slow database look dead, which produced its
slow outliers.

\textbf{LLMs: priors, advice and evidence.} Without guidance, gpt-4o-mini behaves like the runbook: it
restarts first. One paragraph of generic SRE advice did not make the model reason case by case; it
replaced one prior with another (bypass first, in every trial). That helped exactly where restarts are
wrong and hurt where they are right, because a crashed service that times out looks like an unreachable
one in a symptom summary. An explicit investigation step (tools that query readiness, metrics and logs,
and a required hypothesis before deciding) let the advised agent tell the cases apart, making tools plus
advice the most robust LLM configuration, at a cost of seconds per decision. Tools without advice only
added time. No LLM configuration discovered that warming the cache fixes a slow database, a fix the
bandit and the expert both learned from measured outcomes. In this setting, LLMs contributed most as
explainers (post-mortems, questions answered from measured data) and as cautious first responders
behind a shield, and least as optimisers of recovery time; prompt content should be evaluated as
carefully as code, because a single generic paragraph moved recovery times by an order of magnitude in
both directions.

\textbf{Safety is layered.} The shield blocked every unsafe proposal, including repeats proposed by an
LLM agent, and the independent playbook allowlist would have stopped them again. In on-call mode the
human-approval path is fail-safe: silence means no action. The agent's ``confirm before acting'' and
``escalate instead of retrying'' rules prevented two classic automation failures: acting on stale
alerts and restart loops.


\textbf{Resilience as a release criterion.} The seeded bug shows why recovery deserves its own gate:
a remediation path can be syntactically correct, return success at every layer and still do nothing.
Only a test that injects the fault and measures user-visible recovery catches it. The same reasoning
applies to alerts: unit tests for burn-rate rules caught boundary behaviour (a 5\,s blip must not page)
that is otherwise discovered in production. The security stage found real, fixable problems on its
first run (vulnerable build-time dependencies, containers running as root), which illustrates that
DevSecOps gates pay off even in small projects.

\textbf{Threats to validity.} \emph{Internal:} all experiments ran on one laptop; absolute times
include Ansible and probing overhead and would differ on other hardware, but all policies share the
same conditions. Samples are small (two or three trials per cell), so we report ranges rather than
confidence intervals. LLM outputs vary between runs even at temperature~0. The synthetic users and
the on-call agent ran during the ablation but not during the earlier policy comparison, so we compare
LLM variants only within the ablation batch. \emph{Construct:} recovery thresholds (99\%, 300\,ms,
five windows) and the bandit's ``best action'' oracle (highest mean logged reward) are design choices;
other thresholds could change rankings near the boundaries. \emph{External:} the shop has one read
path, seven fault types and no write traffic; real systems have more failure modes, cascading effects
and noisier signals. The expert heuristic was written after seeing the training data, which favours
it on training faults.
